\section{METODOLOGI PENGEMBANGAN}

\subsection{Metodologi yang Dipilih}

Savora dikembangkan menggunakan metodologi \textbf{Scrum} dengan iterasi sprint 2 minggu. Metodologi ini dipilih karena fleksibilitasnya dalam mengakomodasi perubahan kebutuhan dan kemampuannya untuk menghasilkan produk yang selaras dengan ekspektasi pengguna.

\subsection{Alasan Pemilihan Scrum}

\begin{enumerate}[leftmargin=*]
    \item \textbf{Iterasi Cepat} : Memungkinkan pengembangan dan pengujian fitur secara bertahap, sehingga bug dapat teridentifikasi lebih awal.
    \item \textbf{Feedback Berkala} : Stakeholder dapat memberikan masukan setiap akhir sprint untuk perbaikan berkelanjutan.
    \item \textbf{Transparansi} : Daily standup dan sprint review memastikan seluruh tim memiliki pemahaman yang sama tentang progres.
    \item \textbf{Adaptif} : Perubahan prioritas atau kebutuhan dapat diakomodasi tanpa mengganggu seluruh roadmap.
\end{enumerate}

\subsection{Tahapan dan Timeline Pengembangan}

\begin{table}[H]
\centering
\caption{Tahapan Pengembangan Savora}
\label{tab:tahapan}
\begin{tabularx}{\textwidth}{|l|X|X|}
\hline
\textbf{Fase} & \textbf{Aktivitas} & \textbf{Durasi} \\
\hline
Sprint 0 & Setup environment, Wireframe, Database design & 1 minggu \\
\hline
Sprint 1 & Autentikasi, Dashboard Admin & 2 minggu \\
\hline
Sprint 2 & Manajemen Produk UMKM, Marketplace & 2 minggu \\
\hline
Sprint 3 & Food Eligibility, Dynamic Pricing & 2 minggu \\
\hline
Sprint 4 & Checkout, Payment simulation & 2 minggu \\
\hline
Sprint 5 & Tracking, Rating, Notifikasi & 2 minggu \\
\hline
Sprint 6 & Integrasi, Testing, Deployment & 2 minggu \\
\hline
\end{tabularx}
\end{table}

\begin{figure}[H]
\centering
% Placeholder untuk Gantt chart
\fbox{\parbox{0.8\textwidth}{\centering\vspace{3cm}Gantt Chart Pengembangan\\(Akan diganti dengan Gantt chart aktual)\vspace{3cm}}}
\caption{Timeline Pengembangan Savora}
\label{fig:gantt}
\end{figure}

\subsection{Peran Anggota Tim}

\begin{table}[H]
\centering
\caption{Pembagian Tugas Anggota Tim}
\label{tab:peran}
\begin{tabularx}{\textwidth}{|l|Y|Y|}
\hline
\textbf{Role} & \textbf{Anggota} & \textbf{Responsibilitas} \\
\hline
Project Manager & Wa Ode Nur Alia & Koordinasi, perencanaan, komunikasi \\
\hline
Frontend Developer & Richard Firmansyah & Pengembangan UI/UX, React.js \\
\hline
Backend Developer & Muhammad Rifaidi & API, Database, Autentikasi \\
\hline
ML Engineer & Ridwan Hakim Ramadhan & Food Eligibility Score, Dynamic Pricing, Analitik \\
\hline
QA \& DevOps & Nadi Azzada Akbar & Testing, Deployment, CI/CD \\
\hline
\end{tabularx}
\end{table}

\subsection{Tools Pendukung}

\begin{itemize}[leftmargin=*]
    \item \textbf{Project Management} : Notion atau Trello untuk sprint board dan task tracking
    \item \textbf{Communication} : Discord atau WhatsApp Group untuk komunikasi harian
    \item \textbf{Design} : Figma untuk wireframe dan UI design
    \item \textbf{Documentation} : Google Docs untuk dokumentasi teknis
    \item \textbf{Version Control} : Git dengan branch strategy (main, develop, feature branches)
\end{itemize}
